\section{Conclusion}
\label{sec:conclusion}

\modelly{} gives Python modellers the agent-based programming model they know and reduces
what it costs to do research with it. A model keeps its agents as objects and has its
population-wide phases written as array operations, one at a time; on eight workloads those
phases run 2.9 to 33 times faster than per-agent Mesa~3.5 code, the lead growing with the
population on two of the three workloads timed at three sizes, within 1.8 to 4.5 times of hand-written numpy that gives agents up, and counted
populations two orders of magnitude faster; an ablation attributes the gain to the array
form of the phases. Whether the faster forms are the same model is tested: against the author's Mesa~2.3
reference implementations, 29 of 67 statistic--variant comparisons are shown equivalent within half
a standard deviation over 100 paired seeds after Holm correction (63 at nominal significance;
48 shown equivalent within $\pm 5\%$ of Mesa's mean), and on three workloads the whole
curves, the spatial patterns and the times to converge of the faster forms are as close to
Mesa~2.3 as Mesa~3.5's are. What this establishes is close agreement on the dynamics
tested, at the tolerances stated for each quantity, not a general scientific equivalence
across all behaviours: 34 of 52 scalar quantities are shown equivalent after Holm
correction, the rest are undecided at 100 seeds, and one difference is shown, the column
form's epidemic peak 3.3 steps earlier than Mesa~2.3's ($[-6.6, -0.1]$), so a study for
which a peak's timing matters to that precision should set its own tolerance and test it
as \cref{sec:trajectories} does; and Mesa's own examples, migrated with the codemod, give all 35 of
their statistics within the same margin, and the virus example's migrated copy is the same
process as Mesa's, draw for draw, once the random streams are aligned (\cref{sec:samecode}). Around that model, the
steps of a study (sweeps, calibration, sensitivity analysis, emulation, provenance and
archiving) are part of the framework, built on one seeding discipline and checked against
problems whose answers are known, and the reproducibility guarantees are stated by scope:
deterministic on one build; reproducible at ten significant digits on every hosted build
tested, the released wheels installed from PyPI on fifteen (five architectures, three
interpreters) among them; and bit-identical across those fifteen for thirteen of seventeen
examples (twelve of seventeen across the nine builds compiled from the checkout). The same model splits across processes with byte-identical results, runs its array
phases on a GPU (with a bounded gain once construction is counted), connects to live data and
continuous fields, and can be shared as a web page (\cref{tab:scale}); these capabilities are
validated functionality with the performance evidence of \cref{tab:scale}, not general
acceleration claims. For a researcher who
needs many runs, large populations or results that others can check, \modelly{} brings the
tools usually assembled by hand around a simulation library into one package, one seed and one
record. The conditions under which these results hold, and the evidence still to gather, are
set out in \cref{sec:threats}.

\section{Availability}
\label{sec:availability}

\cmd{pip install modelly-abm} installs \modelly{} (Python 3.12 or later), which is imported as
\code{modelly} and is free software under the Apache 2.0 licence. The release workflow builds
compiled wheels for Linux, macOS and Windows (x86-64 and ARM64) on the tag and a pure-Python
wheel for everything else (\cref{tab:determinism}). Source, issue tracker, the tests and the benchmark harness are at
\href{https://github.com/RemiPelloux/Modelly}{\nolinkurl{github.com/RemiPelloux/Modelly}};
documentation, tutorials and the roadmap at
\href{https://remipelloux.github.io/Modelly/}{\nolinkurl{remipelloux.github.io/Modelly}}.
At the tag, the suite runs 3,054 tests, of which 2,996 pass and 58 skip without an optional
library (CuPy, Ray, optuna, xarray, a broker) or a browser, with 97.3\% line coverage (\code{paper/data/test\_suite.json}). The version described here is 1.4.2; the provenance of every measurement (commits,
machine, package versions and the Software Heritage identifier of the archived code) is in
\cref{app:reproduce}, and the data and scripts behind every table and figure of this paper
are in the repository's \code{paper/} folder. To cite \modelly{}, use the metadata in the
repository's \code{CITATION.cff}.

\section*{Declarations}
\addcontentsline{toc}{section}{Declarations}

\paragraph{Software availability.} Name: \modelly{} (\code{modelly-abm} on PyPI). Developer:
R\'emi Pelloux, OpenPro (\href{mailto:rp@openpro.ai}{rp@openpro.ai}). First available: 2026.
Version described: 1.4.2. Language: Python, with five C extension modules. Requirements:
CPython 3.12 or later, \code{numpy}, \code{pandas}, \code{scipy} and \code{tqdm} (and
\code{pythran} on macOS); optional extras for visualisation, geography, networks, experiment
storage, Dask and Ray executors, MPI (\code{mpi4py}), GPUs (CuPy), differentiable models (JAX),
message brokers, an outside optimiser and xarray.
Hardware: any machine that runs CPython; the release workflow builds compiled wheels for
Linux, macOS and Windows on x86-64 and ARM64 and a pure-Python wheel for everything else and
for the browser. Size: 1.17~MB for the
Linux x86-64 wheel and 1.12~MB for the pure one. Licence: Apache 2.0. Source:
\href{https://github.com/RemiPelloux/Modelly}{\nolinkurl{github.com/RemiPelloux/Modelly}}.

\paragraph{Data availability.} The data and scripts behind every measured table and figure are
in the repository's \code{paper/} and \code{benchmarks/} folders (\cref{app:reproduce}).

\paragraph{Funding.} The author received no specific grant for this work.

\paragraph{Competing interests.} The author develops \modelly{}, the software evaluated here.
